\section{Per-sensor attribution}
\label{ch:results-attribution}

Notebook: \file{07\_attribution.ipynb}.  Figure: F5.  Source: \file{results/metrics/attribution.parquet} ($3$ seeds $\times$ $4$ AE-family models $\times$ reconstruction; $3$ seeds $\times$ TranAD $\times$ attention) and \file{results/metrics/attribution\_shap.parquet} ($3$ seeds $\times$ $\{$IF, OCSVM$\}$ $\times$ SHAP) for the two classical baselines that don't have reconstruction-error attribution.

\subsection{Evaluation protocol recap}

For each attack window in HAI test, the model produces a per-feature attribution vector.  We rank features by attribution and compute precision@k --- the fraction of the top-k features that belong to the attacked process (from HAI's \code{attack\_P\{n\}} label).  Random baseline $=$ share of features in the attacked process: $0.48$ for P1, $0.28$ for P2, $0.09$ for P3.

\subsection{Precision@5 (mean $\pm$ std across seeds 7, 42, 123)}

\begin{table}[h]
\centering
\caption{Precision@5 by model and attacked process, averaged over three seeds.  Best per process in bold.}
\label{tab:attr-p5}
\small
\begin{tabular}{lrrr}
\toprule
Model & P1 p@5 & P2 p@5 & P3 p@5 \\
\midrule
Dense AE & $0.69 \pm 0.01$ & $0.74 \pm 0.00$ & $\mathbf{0.33 \pm 0.02}$ \\
LSTM-AE  & $0.74 \pm 0.00$ & $0.09 \pm 0.00$ & $0.20 \pm 0.00$ \\
USAD     & $0.74 \pm 0.00$ & $0.09 \pm 0.00$ & $0.21 \pm 0.02$ \\
TranAD   & $0.58 \pm 0.17$ & $\mathbf{0.35 \pm 0.19}$ & $0.02 \pm 0.01$ \\
Random   & $0.48$ & $0.28$ & $0.09$ \\
\bottomrule
\end{tabular}
\end{table}

\begin{figure}[h]
\centering
\includegraphics[width=0.95\textwidth]{07_attribution/attribution_p_at_5.png}
\caption{Attribution precision@5 by model and attacked process, with $\pm 1\sigma$ error bars across three seeds.  Dashed line is the random baseline for that process.}
\label{fig:attr-p5}
\end{figure}

\subsection{Finding 4.1 --- detection-vs-localization tradeoff (robust across seeds)}

Lift over random baseline for mean p@5:

\begin{table}[h]
\centering
\small
\begin{tabular}{lrrr}
\toprule
Model & P1 lift & P2 lift & P3 lift \\
\midrule
\textbf{Dense AE} & $\mathbf{1.45\times}$ & $\mathbf{2.68\times}$ & $\mathbf{3.72\times}$ \\
LSTM-AE & $1.55\times$ & $0.32\times$ & $2.22\times$ \\
USAD    & $1.55\times$ & $0.33\times$ & $2.33\times$ \\
TranAD  & $1.22\times$ & $1.26\times$ & $0.23\times$ \\
\bottomrule
\end{tabular}
\end{table}

\textbf{Dense AE --- the weakest Phase-2 detector --- is the only AE-family / windowed model whose attribution is reliably above random across all three processes.}  Seed variance is small ($\sim 0.01$--$0.02$ on P3); robust result, not seed artifact.  Finding~4.4 below shows the two classical baselines (IF and OCSVM via SHAP) join Dense AE in this everywhere-above-random group, with OCSVM in particular reaching Dense-AE-level localization despite its bottom-of-the-table detection F1.

The mechanism: Dense AE operates per-timestep, so per-feature reconstruction error is aligned with the current frame's anomaly.  Windowed models reconstruct a $60$-step window and average per-feature error over time; the averaging concentrates mass on features with the largest dynamic-range variation --- HAI's P1 flow and pressure sensors.  A small perturbation to a few discrete P2 flags is lost in the mean.

\textbf{This inverts the detection ranking of Section~\ref{ch:results-detection}}: TranAD is \#1 on HAI eTaPR but near-worst on P3 localization ($0.02$ mean p@5, $0.23\times$ random lift).  For an operator triaging a flagged window, Dense AE's per-feature output is more actionable.  Reporting only detection F1 obscures this.

\subsection{Finding 4.2 --- attention rollout is null (robust across seeds)}
\label{sec:finding-attn-null}

TranAD attention-weighted attribution differs from plain reconstruction attribution by $< 0.01$ on every $(\text{process}, k)$ pair across all three seeds.  The two methods are empirically interchangeable.

Intuition: with a single-layer encoder, rollout $=$ raw attention matrix.  Attention averaged over $60$ output positions on HAI normal data is close to uniform; reweighting timesteps uniformly leaves the per-feature ranking unchanged.  Attention-based XAI for per-sensor localization in ICS windowed AD is not automatically the right tool --- one has to verify attention-over-time is structurally variable before investing.

\subsection{Finding 4.3 --- LSTM-AE $\equiv$ USAD attribution is a genuine structural result}

Across $3$ seeds, LSTM-AE and USAD produce p@k values identical to two decimals on every $(\text{process}, k)$ pair.  Standard deviations are $\leq 0.016$ within each model.  This is not seed noise; the two architectures converge to the same per-feature reconstruction ranking on HAI.

Hypothesis: the ranking is dominated by the MinMax-scaled feature variance structure of HAI's $79$ sensors.  P1 sensors carry the highest dynamic range and feature count, so any architecture with enough capacity to minimise windowed reconstruction error learns to prioritise them.  LSTM-AE's sequence objective and USAD's adversarial dual-head objective both converge to the same optimum in the feature-ranking sense.

If your model class produces the same attribution ranking as a simpler baseline, the extra architectural complexity is buying detection accuracy (USAD and LSTM-AE \emph{do} differ in detection F1 slightly), not attribution quality.  Operators paying the complexity tax for a richer model should audit that they're getting the richer output.

\subsection{Revised framing vs.\ the seed-42-only report}

An earlier draft (commit \texttt{a2957b0}) reported ``TranAD actively anti-localizes on P3 (p@5 $= 0.02$, lift $0.2\times$, worse than flipping coins).''  Multi-seed averaging softens this claim:
\begin{itemize}
  \item P3 remains a TranAD weakness (mean $0.02$, range $[0.02, 0.03]$ across seeds --- consistent, not seed-random).
  \item But TranAD on P2 had large seed variance: p@5 $= 0.23$ at seed $42$ vs.\ $0.57$ at seed $7$.  Mean $0.35$, \emph{above} the $0.28$ random baseline.
\end{itemize}

Revised framing: TranAD \emph{can} localize to P2 with enough seed-trial budget, but its localization is seed-unstable, whereas Dense AE and LSTM-AE/USAD are seed-stable (even if the latter two are structurally stuck on P1).  An argument for multi-seed evaluation in any attribution study.

\subsection{Finding 4.4 --- SHAP fills in the classical baselines}
\label{sec:finding-shap}

The four AE-family / windowed models above use reconstruction-error or attention-rollout attribution, both of which require an internal reconstruction signal.  Isolation Forest and One-Class SVM produce a single anomaly score per row with no per-feature decomposition, so the same evaluation cannot be applied directly.  We close that gap with SHAP: \code{shap.TreeExplainer} on the underlying \code{IsolationForest} ensemble (exact, $\sim$ms/row), and \code{shap.KernelExplainer} on the \code{OneClassSVM} decision function with a deterministic $100$-row background sample drawn from the training-normal split (approximate, $\sim$seconds/row).  The KernelExplainer's predictor wraps \code{-decision\_function} so SHAP explains the project's anomaly-score convention (higher contribution = more anomalous) directly, and the background sample is fixed across the three training seeds so SHAP variance reflects the artifacts, not the sampling.  Evaluation reuses the same \code{precision\_at\_k\_by\_attack} harness from \S\ref{ch:results-attribution}; the only deviation from Phase 4's protocol is that OCSVM evaluation subsamples to $100$ attack rows per process per seed to keep KernelExplainer's runtime bounded --- the precision@k aggregate is robust to this at $n_{\text{attack}} \geq 100$.

\begin{table}[h]
\centering
\caption{SHAP-based precision@$5$ for the two classical baselines, mean over seeds $\{7, 42, 123\}$, alongside the AE-family numbers from Table~\ref{tab:attr-p5}.  Random baselines: $0.48$ for P1, $0.28$ for P2, $0.09$ for P3.}
\label{tab:attr-shap-p5}
\small
\begin{tabular}{llrrr}
\toprule
Model & Method & P1 p@5 & P2 p@5 & P3 p@5 \\
\midrule
Isolation Forest & SHAP-tree   & $0.67 \pm 0.03$ & $0.38 \pm 0.04$ & $0.19 \pm 0.06$ \\
OCSVM            & SHAP-kernel & $0.69 \pm 0.02$ & $\mathbf{0.42 \pm 0.03}$ & $0.25 \pm 0.00$ \\
\midrule
Dense AE & reconstruction & $0.69$ & $0.74$ & $\mathbf{0.33}$ \\
LSTM-AE  & reconstruction & $0.74$ & $0.09$ & $0.20$ \\
USAD     & reconstruction & $0.74$ & $0.09$ & $0.21$ \\
TranAD   & reconstruction & $0.58$ & $0.35$ & $0.02$ \\
\bottomrule
\end{tabular}
\end{table}

\begin{figure}[h]
\centering
\includegraphics[width=0.95\textwidth]{07_attribution/shap_if_ocsvm_pak.png}
\caption{Per-sensor attribution precision@$5$ across all six models on HAI test, with $\pm 1\sigma$ error bars across three seeds.  Classical baselines (IF, OCSVM) use SHAP; AE-family models use reconstruction-error attribution from \S\ref{ch:results-attribution}.  Random baseline shown as a dashed line per panel.}
\label{fig:attr-shap}
\end{figure}

Both classical baselines clear the random baseline on every process at p@$5$: IF lifts of $1.39\times$ / $1.35\times$ / $2.18\times$ on P1 / P2 / P3, OCSVM lifts of $1.44\times$ / $1.50\times$ / $2.79\times$.  The interesting result is OCSVM: it ties Dense AE on P1 ($0.69$ vs.\ $0.69$), beats every windowed model on P2 ($0.42$ vs.\ $0.09$ for LSTM-AE / USAD and $0.35$ for TranAD), and sits between LSTM-AE / USAD and Dense AE on P3.  Phase~2 ranks OCSVM near the bottom of the HAI detection table (eTaPR $0.170$, fifth of six), so this is the same detection-vs-localization inversion that Dense AE exhibited in Finding 4.1, now appearing in a kernel method too: the two models with the weakest absolute detection on HAI produce the most actionable per-feature attributions across all three attacked processes.

The mechanism for OCSVM's localization quality is plausible: \code{shap.KernelExplainer} on the OCSVM decision function attributes per-feature contribution by perturbing each feature against the $100$-row training-normal background; features whose perturbation moves the kernel-distance decision the most are exactly those that distinguish the test row from the normal manifold the SVM learned.  That signal is per-row and unaveraged over time, like Dense AE's, so it preserves discrete-flag attacks (P2) that the windowed models smear away.  IF's \code{TreeExplainer} attribution is weaker on P2 / P3 because the model itself isolates rows with fewer effective splits when many features look normal; on P1 (sustained large-magnitude attacks across many sensors) IF localizes well, but on the discrete-flag and rare-process attacks it has less to attribute.

\textbf{Caveat.}  KernelExplainer with a $100$-row background is approximate: the qualitative ranking is robust to background size in our setup, but the exact contribution magnitudes are not, and OCSVM's $100$ attack rows per process per seed is enough for ranking but tighter than the AE-family's full attack set ($n_{\text{P1}} \approx 7000$, $n_{\text{P2}} \approx 1300$, $n_{\text{P3}} \approx 410$ per seed).  As with the AE-family results above, p@k is computed at the process granularity HAI provides; per-sensor evaluation would need labelled per-sensor ground truth which HAI 21.03 does not release.

\subsection{Limitations}

\begin{enumerate}
  \item \textbf{Process-level eval is weaker than per-sensor eval.}  HAI 22.04's README has more detailed per-attack descriptions; a hand-curated per-sensor ground truth on 22.04 would strengthen the claim.
  \item \textbf{Dense AE is non-windowed; LSTM-AE, USAD, TranAD are windowed.}  The localization-vs-detection tradeoff is partly confounded with windowing.  A windowed MLP AE would disentangle; out of scope here.
  \item \textbf{Attention rollout is a single-layer approximation.}  A deeper TranAD could give richer signals; deeper TranAD also takes hours per run on our 4\,GB GPU.
  \item \textbf{KernelExplainer is approximate.}  OCSVM SHAP values are computed with a $100$-row background sample and $128$ Monte-Carlo perturbations per row; tighter settings exceeded our compute budget.  Qualitative rankings should be stable but exact contribution magnitudes are not (\S\ref{sec:finding-shap}).
\end{enumerate}
